\section{Cách xác định output của một bộ PID}

Một bộ PID không tự quyết định output là gia tốc, lực, moment hay PWM.
Người thiết kế quyết định output dựa trên câu hỏi: "Tầng sau của PID cần nhận đại lượng gì?"

Có thể viết nguyên tắc này dưới dạng:

\begin{equation}
    \boxed{
    \text{Đơn vị output của PID}
    =
    \text{Đơn vị input mà tầng sau cần}
    }
\end{equation}

Ví dụ, nếu tầng sau cần gia tốc:

\begin{equation}
    u = a
\end{equation}

Nếu tầng sau cần lực:

\begin{equation}
    u = F
\end{equation}

Nếu tầng sau cần moment:

\begin{equation}
    u = \tau
\end{equation}

Nếu tầng sau cần PWM:

\begin{equation}
    u = PWM
\end{equation}

\subsection{Output của các tầng PID trong project}

\subsubsection{Position PID}

Tầng Position PID nhận sai số vị trí và biến nó thành lệnh gia tốc/góc nghiêng mong muốn \cite{mahony2012multirotor}:

\begin{equation}
    e_x,\quad e_y
\end{equation}

và tạo gia tốc ngang mong muốn:

\begin{equation}
    a_x,\quad a_y
\end{equation}

Sau đó đổi thành góc nghiêng mong muốn:

\begin{equation}
    a_x,\ a_y
    \longrightarrow
    \phi_d,\ \theta_d
\end{equation}

Vì vậy:

\begin{equation}
    \boxed{
    PositionPID
    \longrightarrow
    roll\_target,\ pitch\_target
    }
\end{equation}

\subsubsection{Altitude PID}

Tầng Altitude PID điều khiển độ cao $z$.

Sai số độ cao:

\begin{equation}
    e_z = z_d - z
\end{equation}

Sai số vận tốc đứng:

\begin{equation}
    e_{v_z} = v_{zd} - v_z
\end{equation}

Trong project, output của tầng này là PWM nền:

\begin{equation}
    PWM_{base}
    =
    PWM_{hover}
    +
    K_{p_z}e_z
    +
    K_{d_z}e_{v_z}
    +
    K_{i_z}\int e_z\,dt
\end{equation}

Vì vậy:

\begin{equation}
    \boxed{
    AltitudePID
    \longrightarrow
    PWM_{base}
    }
\end{equation}

\subsubsection{Attitude PID}

Tầng Attitude PID nhận góc mong muốn:

\begin{equation}
    \phi_d,\quad \theta_d
\end{equation}

và góc hiện tại:

\begin{equation}
    \phi,\quad \theta
\end{equation}

Sai số roll:

\begin{equation}
    e_\phi = \phi_d - \phi
\end{equation}

Sai số pitch:

\begin{equation}
    e_\theta = \theta_d - \theta
\end{equation}

Trong project, output của tầng này là PWM correction:

\begin{equation}
    \Delta PWM_\phi
    =
    K_{p_\phi}e_\phi
    +
    K_{i_\phi}\int e_\phi\,dt
    -
    K_{d_\phi}p
\end{equation}

\begin{equation}
    \Delta PWM_\theta
    =
    K_{p_\theta}e_\theta
    +
    K_{i_\theta}\int e_\theta\,dt
    -
    K_{d_\theta}q
\end{equation}

Trong đó:

\begin{equation}
    p = \dot{\phi}
\end{equation}

\begin{equation}
    q = \dot{\theta}
\end{equation}

Vì vậy:

\begin{equation}
    \boxed{
    AttitudePID
    \longrightarrow
    \Delta PWM_\phi,\ \Delta PWM_\theta
    }
\end{equation}

\subsubsection{Yaw PID}

Sai số yaw:

\begin{equation}
    e_\psi = wrap(\psi_d - \psi)
\end{equation}

Output của yaw PID là PWM correction theo yaw:

\begin{equation}
    \Delta PWM_\psi
    =
    K_{p_\psi}e_\psi
    +
    K_{i_\psi}\int e_\psi\,dt
    -
    K_{d_\psi}r
\end{equation}

Trong đó:

\begin{equation}
    r = \dot{\psi}
\end{equation}

Vì vậy:

\begin{equation}
    \boxed{
    YawPID
    \longrightarrow
    \Delta PWM_\psi
    }
\end{equation}

\subsection{Mixer}

Mixer nhận các đại lượng:

\begin{equation}
    PWM_{base},\quad
    \Delta PWM_\phi,\quad
    \Delta PWM_\theta,\quad
    \Delta PWM_\psi
\end{equation}

và tạo PWM cho từng motor:

\begin{equation}
    PWM_1
    =
    PWM_{base}
    +
    \Delta PWM_\phi
    -
    \Delta PWM_\theta
    -
    \Delta PWM_\psi
\end{equation}

\begin{equation}
    PWM_2
    =
    PWM_{base}
    -
    \Delta PWM_\phi
    -
    \Delta PWM_\theta
    +
    \Delta PWM_\psi
\end{equation}

\begin{equation}
    PWM_3
    =
    PWM_{base}
    -
    \Delta PWM_\phi
    +
    \Delta PWM_\theta
    -
    \Delta PWM_\psi
\end{equation}

\begin{equation}
    PWM_4
    =
    PWM_{base}
    +
    \Delta PWM_\phi
    +
    \Delta PWM_\theta
    +
    \Delta PWM_\psi
\end{equation}

Vì vậy:

\begin{equation}
    \boxed{
    Mixer
    \longrightarrow
    PWM_1,\ PWM_2,\ PWM_3,\ PWM_4
    }
\end{equation}

\subsection{Tóm tắt nguyên tắc chọn output PID}

Nguyên tắc chung:

\begin{equation}
    \boxed{
    \text{PID output phụ thuộc vào tầng sau cần gì}
    }
\end{equation}

Trong project hiện tại:

\begin{equation}
    PositionPID
    \longrightarrow
    a_x,\ a_y
    \longrightarrow
    \phi_d,\ \theta_d
\end{equation}

\begin{equation}
    AltitudePID
    \longrightarrow
    PWM_{base}
\end{equation}

\begin{equation}
    AttitudePID
    \longrightarrow
    \Delta PWM_\phi,\ \Delta PWM_\theta
\end{equation}

\begin{equation}
    YawPID
    \longrightarrow
    \Delta PWM_\psi
\end{equation}

\begin{equation}
    Mixer
    \longrightarrow
    PWM_1,\ PWM_2,\ PWM_3,\ PWM_4
\end{equation}

\begin{equation}
    MotorModel
    \longrightarrow
    T_1,\ T_2,\ T_3,\ T_4
\end{equation}

Do đó, trong project này:
\begin{itemize}
    \item Position PID tạo gia tốc rồi đổi thành góc nghiêng mong muốn.
    \item Altitude PID tạo PWM nền.
    \item Attitude PID tạo PWM hiệu chỉnh roll/pitch.
    \item Yaw PID tạo PWM hiệu chỉnh yaw.
    \item Mixer tạo PWM cho từng motor.
    \item Motor model biến PWM thành lực nâng.
\end{itemize}
